\subsection{Jaka jest różnica między przestrzenią Hilberta a przestrzenią fizyczną?}

odpowiedź

\subsection{Czym jest operator liniowy w mechanice kwantowej?}

odpowiedź

\subsection{Zdefiniuj operator hermitowski.}

odpowiedź

\subsection{Jakie wielkości fizyczne odpowiadają operatorom hermitowskim?}

odpowiedź

\subsection{Czym są wartości własne i wektory własne w kontekście pomiarów kwantowych?}

odpowiedź

\subsection{Jakie jest fizyczne znaczenie iloczynu skalarnego w przestrzeni Hilberta?}

odpowiedź

\subsection{Wyjaśnij znaczenie normalizacji stanów kwantowych.}

odpowiedź

\subsection{Co oznacza, że dwa stany kwantowe są ortogonalne?}

odpowiedź

\subsection{Wyjaśnij pojęcie superpozycji stanów.}

odpowiedź

\subsection{Jakie są relacje komutacyjne operatorów położenia i pędu?}

odpowiedź

\subsection{Co oznacza, że dwa operatory komutują?}

odpowiedź

\subsection{Czym jest pełny układ bazowy w przestrzeni Hilberta?}

odpowiedź

\subsection{Co to jest operator unitarny?}

odpowiedź

\subsection{Jak transformacje unitarne zachowują prawdopodobieństwo?}

odpowiedź

\subsection{Wyjaśnij ewolucję czasową za pomocą operatora ewolucji.}

odpowiedź

\subsection{Wyjaśnij pojęcie stanów stacjonarnych.}

odpowiedź

\subsection{Czym jest pakiet falowy?}

odpowiedź

\subsection{Co oznacza transformacja Fouriera funkcji falowej i jakie ma znaczenie fizyczne?}

odpowiedź

\subsection{Wyjaśnij zjawisko tunelowania kwantowego.}

odpowiedź

\subsection{Wyjaśnij transformację Fouriera między przestrzenią położeń a pędu.}

odpowiedź

\subsection{Czym są stany swobodne i związane?}

odpowiedź

\subsection{Jakie są warunki ciągłości dla funkcji falowej i jej pochodnych?}

odpowiedź

\subsection{Co to jest operator momentu pędu?}

odpowiedź

\subsection{Jakie są relacje komutacyjne dla składników momentu pędu?}

odpowiedź

\subsection{Jakie są wartości własne $L^2$ i $L_z$?}

odpowiedź

\subsection{Jaki jest związek między funkcjami $Y_{lm}$ a momentem pędu?}

odpowiedź

\subsection{Co to jest obrót w przestrzeni i jak działa na funkcje falowe?}

odpowiedź

\subsection{Co to jest moment własny (spin)?}

odpowiedź

\subsection{Jakie są wartości dozwolone dla spinu cząstki?}

odpowiedź

\subsection{Czym są macierze Pauliego?}

odpowiedź

\subsection{Zdefiniuj operator spinowy $S$.}

Operator spinowy $S$ to operator kwantowomechaniczny, który opisuje wewnętrzny moment pędu cząstki (np. elektronu).


\subsection{Co to jest doświadczenie Stern-Gerlacha i co pokazuje?}

Doświadczenie to polegało na przepuszczeniu wiązki atomów srebra przez niejednorodne pole magnetyczne i obserwacji obrazu wiązki na ekranie (np. kliszy). Wyniki pokazały, że moment pędu (dokładnie spin) przyjmuje wartości skwantowane — strumień się rozdziela na dwie wiązki, co odpowiada dwu stanom spinu (+1/2 lub -1/2).

\subsection{Co to znaczy, że spin nie ma klasycznego odpowiednika?}

Spin to wewnętrzna własność cząstek, która nie znajduje odzwierciedlenia w jakimś ruchu obrotowym czy obrocie ładunku, tak jak moment pędu orbitalnego. Nie można sobie go wyobrazić jak kręcącą się kulkę — to wielkość czysto kwantowa.

\subsection{Jakie są stany własne atomu wodoru?}

odpowiedź

\subsection{Jakie są liczby kwantowe w atomie wodoru?}
\begin{itemize}
  \item $n$ -- główna liczba kwantowa (przyjmuej wartości $n = 1,2,3,\dots$,
  \item $l$ -- orbitalna liczba kwanntowa ($l = 0,1,2,\dots, n-1$,
  \item $m_l$ -- magnetyczna liczba kwantowa ($-l,-(l-1),\dots, l-1, l$)

\subsection{Wyjaśnij degenerację poziomów energetycznych w atomie wodoru.}

sytuacja, gdy kilka różnych stanów kwantowych ma tę samą energię. \\
W przypadku atomu wodoru, stany o tym samym numerze powłoki głównej (n), ale różniących się liczbami kwantowymi orbitalnymi (l) i magnetycznymi (ml), mogą mieć tę samą energię w idealnym modelu atomu.

\subsection{Na czym polega zakaz Pauliego?}

odpowiedź

\subsection{Czym jest symetria permutacji fermionów i bozonów?}

odpowiedź

\subsection{Czym jest model Bohra?}

odpowiedź

\subsection{Czym jest energia jonizacji?}

odpowiedź

\subsection{Czym jest efekt fotoelektryczny?}

odpowiedź
